\begin{table}[H]
\small
\begin{tabular}{@{}lllcccccccc@{}}
\toprule
\textbf{Metric} & \textbf{Space} & \textbf{Interp} & \textbf{Grad Max} & \textbf{Grad Avg} & \textbf{Sat Avg} & \textbf{Int Avg} & \textbf{\#Colors} & \textbf{Eff. Palette} & \textbf{Hue H (bits)} & \textbf{Mean NN $\Delta E$}\\
\midrule
Euclidean & sRGB (cube corners) & sRGB & 48.00 & 2.988 & 0.765 & 0.895 & 4913 & 11.3 & 3.26 & 3.65\\
Euclidean & CIELAB & sRGB & 45.66 & 3.461 & 0.928 & 0.892 & 100166 & 4.1 & 3.38 & 1.14\\
Euclidean & CIELAB ($\Delta E_{76}$ = this row) & CIELAB & 101.19 & 3.430 & 0.932 & 0.894 & 66378 & 3.9 & 3.17 & 1.12\\
Euclidean & OKLAB & sRGB & 72.75 & 2.674 & 0.251 & 0.270 & 42808 & 2.0 & 1.21 & 1.07\\
\bottomrule
\end{tabular}
\caption{Color-space and interpolation-space comparison. All rows fix the argmax metric (Euclidean) and the geometry (no smoothing); \emph{Space} is the intermediate color space used for the farthest-color search, and \emph{Interp} is where the output LUT interpolates between chosen hull vertices. The top row is the trivial sRGB baseline (the hull of the input 0-255 cube is its 8 corners, so each input picks the opposite corner). CIELAB + CIELAB-interp is Kwon 2019's $\Delta E_{76}$ baseline.}
\label{tab:space_and_interp}
\end{table}
